\documentclass[10pt,a4paper]{amsart}
\usepackage[margin=19mm]{geometry}
\usepackage{amsmath,amssymb,mathtools}
\usepackage[scr=rsfs,cal=euler]{mathalfa}
\usepackage{xfrac}
\usepackage[unicode,hidelinks,bookmarks=false]{hyperref}
\newcommand{\ie}{i.e.\ }
\newcommand{\cf}{cf.\ }
\newcommand{\resp}{respectively\ }
\newcommand{\Subsection}{\subsection}
\providecommand{\nobreakdash}{\nobreak\hbox{-}}
\setlength{\emergencystretch}{2em}
\multlinegap0pt
\usepackage{mathtools}
\usepackage{tikz}
\usepackage{dsfont}
\newtheorem{theorem}{Theorem}
\newcommand{\Cons}[2]{#1 \ConsSym #2}
\newcommand{\ConsI}[2]{#1 \ConsSym^* #2}
\newcommand{\ConsSym}{::}
\newcommand{\List}[1]{\mathsf{List}(#1)}
\newcommand{\Nil}{\mathsf{nil}}
\newcommand{\NilI}{\mathsf{nil}^*}
\newcommand{\Rec}[4]{\mathsf{rec}_{#1}(#2, #3, #4)}
\newcommand{\projl}[1]{\pi_1 \> #1}
\title{Impredicativity in Linear Dependent Type Theory}
\author{Sam Speight, Niels van der Weide}
\date{Source: arXiv:2602.08846v1 (CC BY 4.0)}
\begin{document}
\maketitle

\noindent\emph{Source and licence.} Impredicativity in Linear Dependent Type Theory. Version arXiv:2602.08846v1 (CC BY 4.0), distributed by arXiv under the Creative Commons Attribution 4.0 International licence, http://creativecommons.org/licenses/by/4.0/, as recorded in the arXiv licence metadata for this exact version and shown on the abstract page https://arxiv.org/abs/2602.08846v1. Full licence notice: \texttt{licenses/arxiv-2602.08846-CC-BY-4.0.txt}.\par\smallskip

\noindent\textit{Editorial scope.} Counted unit: the article's theorem thm:initial-impred (source0.tex lines 1571--1574). The article's complete original proof of it is delivered with it (source0.tex lines 1576--1613, one \texttt{proof} environment of 38 lines). No mathematics is written, repaired or re-derived: the statement and the proof are the article's own lines, in the article's own notation, and no retained line is substituted or re-flowed.  No further source-local context is needed: every cross-reference of every retained line resolves inside the two delivered spans.  Nothing was omitted from any retained span, so the package carries no omission record.  No retained line contains a citation, so the wrapper carries no bibliography and no bibliography entry.  No retained line contains a figure environment, an image-inclusion command or drawing code, so no image is delivered and the package declares no asset dependency.  Editorial formatting only, in addition: an AMS \texttt{amsart} preamble that restates the article's own declarations and macros used by the retained lines, namely the environments \texttt{theorem} on separate counters, together with the standard packages those lines need.  The retained environments print in this document as Theorem 1 in order of appearance, whereas the article prints the same items with the numbering of its own full document.  The complete native source of the article as distributed in the arXiv e-print for this exact version is bundled unchanged as \texttt{sources/194277/source0.tex}.
\begin{theorem}
\label{thm:initial-impred}
The algebra $\List{A}$ is initial.
\end{theorem}

\begin{proof}
Given an algebra $X$,
we define $f$ to be $\lambda (l : \List{A}) . \Rec{X}{n_X}{c_X}{l}$.
%We first show that for each algebra $X$
%we have a morphism $f : \List{A} \to X$.
%The map $f$ sends $l : \List{A}$ to $\Rec{X}{n_X}{c_X}{\projl{l}}$.
Note that
\begin{align*}
&\Rec{X}{n_X}{c_X}{\projl{\Nil}} = \Rec{X}{n_X}{c_X}{\NilI} = n_X
\\
&\Rec{X}{n_X}{c_X}{\projl{(\Cons{a}{l})}} = \Rec{X}{n_X}{c_X}{\ConsI{a}{l}} = c_X \> a \> (\Rec{X}{n_X}{c_X}{l})
\end{align*}

To prove that the map $f$ is unique,
it suffices to show that $\Rec{\List{A}}{\Nil}{\ConsSym}{l} = l$
for each $l : \List{A}$.
By function extensionality,
we need to show for all $X$, $n_X$, and $c_X$ that
\[
\Rec{\List{A}}{\Nil}{\ConsSym}{l} \> X \> n_X \> c_X = l \> X \> n_X \> c_X.
\]
By definition,
we have $l \> X \> n_X \> c_X = \Rec{X}{n_X}{c_X}{l}$
and that
\[
\Rec{\List{A}}{\Nil}{\ConsSym}{l} \> X \> n_X \> c_X
=
\Rec{X}{n_X}{c_X}{\Rec{\List{A}}{\Nil}{\ConsSym}{l}}.
\]
It suffices to show that
\[
\Rec{X}{n_X}{c_X}{\Rec{\List{A}}{\Nil}{\ConsSym}{l}} 
=
\Rec{X}{n_X}{c_X}{l}.
\]
Since $\lambda l. \Rec{X}{n_X}{c_X}{l}$ is a morphism,
the theorem follows from the fact that $l : \List{A}$.
\end{proof}

\end{document}
